\documentclass[11pt]{amsart}

\usepackage[utf8]{inputenc}
\usepackage[T1]{fontenc}
\usepackage{amsmath,amssymb,amsthm,mathtools}
\usepackage[margin=2.4cm,top=2.5cm,bottom=2.5cm]{geometry}
\usepackage{enumitem}
\usepackage{xcolor}
\usepackage{hyperref}

\definecolor{seminarblue}{RGB}{20,40,80}
\hypersetup{colorlinks=true, linkcolor=seminarblue, citecolor=seminarblue,
  urlcolor=seminarblue,
  pdftitle={From spectra to univalence: a first tour of HoTT},
  pdfauthor={C. Galindo}}

\setlength{\parindent}{0pt}
\setlength{\parskip}{0.32em}

\theoremstyle{plain}
\newtheorem{theorem}{Theorem}[section]
\newtheorem{proposition}[theorem]{Proposition}
\newtheorem{lemma}[theorem]{Lemma}
\newtheorem{corollary}[theorem]{Corollary}
\newtheorem{axiom}[theorem]{Axiom}

\theoremstyle{definition}
\newtheorem{definition}[theorem]{Definition}
\newtheorem{example}[theorem]{Example}

\theoremstyle{remark}
\newtheorem{remark}[theorem]{Remark}

\newcommand{\ZZ}{\mathbb{Z}}
\newcommand{\Uni}{\mathcal{U}}
\newcommand{\isProp}{\mathsf{isProp}}
\newcommand{\isContr}{\mathsf{isContr}}
\newcommand{\refl}{\mathsf{refl}}
\newcommand{\Scirc}{S^1}

\begin{document}

\begin{center}
{\Large\bfseries From spectra to univalence: a first tour of homotopy type theory}\\[1em]
{\large César Galindo}\\[0.3em]
{\small Expository note --- written for a reader who already thinks in
terms of spaces, homotopies, and fibrations (as one does after a thesis
on orthogonal spectra), meeting type theory for the first time.}
\end{center}

\vspace{0.5em}

\begin{center}
\begin{minipage}{0.86\textwidth}
\small
\textsc{Aim.} To build up the basic notions of Homotopy Type Theory ---
identity types, the univalence axiom, mere propositions, and higher
inductive types --- entirely as translations of things a homotopy
theorist already knows: paths, homotopies of maps, contractibility, and
a CW-type presentation of the circle. The payoff, exactly as in
\cite[\S 8.1]{HoTTBook}, is a synthetic computation of $\pi_1(\Scirc)$
that never mentions points, open sets, or a model of space at all.
\end{minipage}
\end{center}

\vspace{1em}

\section{Types as spaces}\label{sec:types-as-spaces}

The guiding dictionary of \cite[\S 2.1]{HoTTBook} is: a \emph{type} $A$
behaves like a space, a \emph{term} $a:A$ behaves like a point of that
space, and a \emph{dependent type} $B(x)$ for $x:A$ behaves like a
fibration $B\to A$. Under this dictionary, equality itself becomes
homotopical: for $x,y:A$ there is not a proposition but a \emph{type}
$x=_A y$, whose terms are witnesses that $x$ and $y$ agree --- to be
read as \emph{paths} from $x$ to $y$ in the space $A$.

\begin{definition}[Identity type, {\cite[\S 1.12]{HoTTBook}}]\label{def:identity-type}
  For a type $A$ and $x,y:A$, the \emph{identity type} $x=_A y$ is the
  type of paths from $x$ to $y$. Every $x:A$ has a canonical path
  $\refl_x : x=_A x$.
\end{definition}

\begin{theorem}[Path induction, {\cite[\S 1.12]{HoTTBook}}]\label{thm:path-induction}
  Let $A$ be a type and let
  \[
    D : \prod_{x,y:A} (x=_A y) \to \Uni
  \]
  be a family of types indexed by pairs of points of $A$ together with a
  path between them. If there is a term $d(x) : D(x,x,\refl_x)$ for
  every $x:A$, then there is a term $f(x,y,p):D(x,y,p)$ for every
  $x,y:A$ and every $p:x=_A y$, with $f(x,x,\refl_x)\equiv d(x)$.
\end{theorem}

Path induction says exactly that, to construct something along an
arbitrary path, it suffices to construct it along the constant path.
Every basic homotopy-theoretic fact about identity types --- symmetry,
transitivity, functoriality of paths under maps --- is proved this way,
mirroring how one proves facts about a space by proving them on a
contractible neighborhood of the diagonal.

\section{Function extensionality and the univalence axiom}\label{sec:univalence}

Two functions that agree pointwise are, in ordinary mathematics, simply
equal; in type theory this has to be an axiom, and it turns out to be a
special case of something much stronger.

\begin{definition}[Equivalence, {\cite[\S 4.1, \S 4.4]{HoTTBook}}]\label{def:equivalence}
  A function $f:A\to B$ is an \emph{equivalence} if its homotopy fibers
  $\mathrm{fib}_f(b) := \sum_{a:A}(f(a)=_B b)$ are contractible for
  every $b:B$. Write $A\simeq B$ for the type of equivalences between
  $A$ and $B$.
\end{definition}

For every type $A$ there is a canonical equivalence $\mathrm{id}_A:A\simeq A$, hence a canonical map
\[
  \mathsf{idtoeqv} \;:\; (A=_{\Uni} B) \longrightarrow (A\simeq B)
\]
sending $\refl_A$ to $\mathrm{id}_A$, for any two types $A,B$ in a
universe $\Uni$.

\begin{axiom}[Univalence, {\cite[\S 2.10]{HoTTBook}}]\label{ax:univalence}
  For every universe $\Uni$ and all $A,B:\Uni$, the map
  $\mathsf{idtoeqv}:(A=_{\Uni}B)\to(A\simeq B)$ is itself an equivalence.
\end{axiom}

Informally: \emph{equivalent types are equal}, and equal in a way that
is itself faithfully recorded by the equivalence one started with. This
is Voevodsky's univalence axiom, the one that gives this repository its
name and this document's subtitle: it is the single extra principle
that turns bare type theory into a foundation where isomorphism
invariance is built in at the syntactic level, rather than imposed
afterwards by fiat.

\begin{theorem}[{\cite[\S 2.9, \S 4.9]{HoTTBook}}]\label{thm:univalence-implies-funext}
  Univalence (Axiom~\ref{ax:univalence}) implies function extensionality:
  for $f,g:\prod_{x:A}B(x)$, the canonical map
  $(f=g)\to\prod_{x:A}(f(x)=g(x))$ is an equivalence.
\end{theorem}

So the pointwise-equality principle a working mathematician reaches for
instinctively is not a separate axiom bolted onto type theory: it falls
out of univalence, the same way many closure properties of weak
equivalences in a model category fall out of one saturation axiom
rather than needing to be checked case by case.

\section{Mere propositions}\label{sec:mere-props}

Not every type behaves like a genuine ``space with interesting paths'';
some behave like ordinary logical propositions, with at most one path
between any two of their points.

\begin{definition}[{\cite[\S 3.3, \S 3.11]{HoTTBook}}]\label{def:mere-prop}
  A type $A$ is a \emph{mere proposition}, $\isProp(A)$, when
  $\prod_{x,y:A}(x=_A y)$ is inhabited. It is \emph{contractible},
  $\isContr(A)$, when $\sum_{a:A}\prod_{x:A}(a=x)$ is inhabited: a mere
  proposition that is, in addition, inhabited.
\end{definition}

\begin{remark}\label{rmk:contractible-fibers}
  Definition~\ref{def:equivalence} can now be read literally: $f$ is an
  equivalence exactly when each of its fibers is a contractible type,
  i.e.\ a space that is homotopy equivalent to a point --- precisely
  the fiberwise condition a homotopy theorist would write down for a
  weak equivalence of spaces or of orthogonal spectra.
\end{remark}

\section{A higher inductive type: the circle}\label{sec:circle}

Ordinary inductive types are generated by point constructors. A
\emph{higher} inductive type is also allowed path constructors ---
generators of its identity type --- which is exactly how one builds
genuinely homotopical spaces synthetically.

\begin{definition}[The circle, {\cite[\S 6.3--6.4]{HoTTBook}}]\label{def:circle}
  The type $\Scirc$ is generated by
  \begin{itemize}[leftmargin=1.6em]
    \item a point $\mathsf{base}:\Scirc$, and
    \item a path $\mathsf{loop}:\mathsf{base}=_{\Scirc}\mathsf{base}$,
  \end{itemize}
  together with an induction principle: to construct a section of a
  family $D:\Scirc\to\Uni$, it suffices to give a point
  $d_{\mathsf{base}}:D(\mathsf{base})$ together with a path
  $d_{\mathsf{base}}=d_{\mathsf{base}}$ lying over $\mathsf{loop}$.
\end{definition}

No quotient of an interval, no CW attaching map, no point-set model is
needed to say this: the circle is presented purely by what one is
allowed to induct on.

\section{Payoff: \texorpdfstring{$\pi_1(\Scirc)\simeq\ZZ$}{pi\_1(S1) = Z} synthetically}\label{sec:payoff}

\begin{theorem}[{\cite[Thm.~8.1.1]{HoTTBook}}]\label{thm:pi1-circle}
  $(\mathsf{base}=_{\Scirc}\mathsf{base}) \;\simeq\; \ZZ$, and hence
  $\pi_1(\Scirc,\mathsf{base})\cong\ZZ$ as groups.
\end{theorem}

\begin{remark}[Sketch of the encode--decode method, {\cite[\S 8.1]{HoTTBook}}]\label{rmk:encode-decode}
  The proof constructs a family $\mathsf{code}:\Scirc\to\Uni$ by the
  recursion principle of Definition~\ref{def:circle}, sending $\mathsf{base}$ to
  $\ZZ$ and transporting along $\mathsf{loop}$ via the
  successor equivalence $\ZZ\simeq\ZZ$; univalence
  (Axiom~\ref{ax:univalence}) is exactly what makes this transport along a
  \emph{loop} land back on $\ZZ$ rather than on some new, unnamed type.
  Maps $\mathsf{encode}:(\mathsf{base}=\mathsf{base})\to\mathsf{code}(\mathsf{base})$
  and $\mathsf{decode}$ in the other direction are then shown to be
  mutually inverse, by path induction (Theorem~\ref{thm:path-induction}) on one
  side and induction on $\ZZ$ on the other. Full details, including the
  fiberwise statement over all of $\Scirc$ that makes the induction go
  through, are in \cite[\S 8.1]{HoTTBook}.
\end{remark}

\begin{remark}[Why this belongs next to a thesis on orthogonal spectra]\label{rmk:bridge-to-thesis}
  Computing $\pi_1(\Scirc)$ is, of course, elementary in classical
  homotopy theory. What is new here is not the answer but the method:
  the computation above never leaves type theory to invoke a
  topological model of the circle, a simplicial set, or a spectrum ---
  it is carried out entirely by induction principles internal to the
  type $\Scirc$ itself. The spectral-sequence techniques of
  \texttt{Theses\_Galindo.tex} compute homotopy invariants \emph{of} a
  fixed point-set or simplicial model; univalent foundations instead
  asks what invariants can be defined \emph{without ever choosing a
  model at all}. Whether, or how, the orthogonal tower of that thesis
  admits a synthetic, univalent presentation is not addressed here ---
  it is left as an open question, not a claim.
\end{remark}

\section{Takeaways}\label{sec:takeaways}

\begin{itemize}[leftmargin=1.6em]
  \item Equality between points of a type is itself a type, built by
    the identity-type former, and everything about it is proved by
    path induction (Definition~\ref{def:identity-type} and Theorem~\ref{thm:path-induction}).
  \item Univalence says equivalent types are equal, and it alone forces
    function extensionality as a theorem rather than a separate
    postulate (Axiom~\ref{ax:univalence} and Theorem~\ref{thm:univalence-implies-funext}).
  \item Mere propositions and contractible types
    (Definition~\ref{def:mere-prop}) let Definition~\ref{def:equivalence}'s
    fiberwise definition of equivalence read exactly like a homotopy
    theorist's definition of weak equivalence.
  \item Higher inductive types generate spaces by their induction
    principle alone; the circle (Definition~\ref{def:circle}) is the simplest
    nontrivial example.
  \item The synthetic computation $\pi_1(\Scirc)\simeq\ZZ$
    (Theorem~\ref{thm:pi1-circle}) is the standard first demonstration that
    this machinery computes real homotopy-theoretic invariants without
    a background model of space.
\end{itemize}

\begin{thebibliography}{9}

\bibitem[UFP13]{HoTTBook}
The Univalent Foundations Program,
\emph{Homotopy Type Theory: Univalent Foundations of Mathematics},
Institute for Advanced Study, 2013.
Book version \texttt{first-edition-697-gcffb6ea}.
\url{http://homotopytypetheory.org/book/}.

\end{thebibliography}

\end{document}
